% ECONOMICS_PROBLEM: {"id": "C65-26", "title": "Quadratic BSDE uniqueness with only one exponential moment", "jel_code": "C65", "source_id": "OP-0565"}
% FIRST_POSTED: 2026-10-09
% ATTEMPT_STATUS: unresolved
% ENTRY_KIND: research_attempt
% !TeX program = pdflatex
% Self-contained: no external input or bibliography files required.
\documentclass[11pt,a4paper]{article}
\usepackage[T1]{fontenc}
\usepackage[utf8]{inputenc}
\usepackage{lmodern}
\usepackage[margin=24mm]{geometry}
\usepackage{amsmath,amssymb,amsthm,mathtools}
\usepackage{booktabs,longtable,array,enumitem,microtype,xurl,hyperref}
\hypersetup{colorlinks=true,linkcolor=blue,citecolor=blue,urlcolor=blue,pdftitle={Checked research attempts: nine economics and stochastic-analysis problems}}
\setlength{\parindent}{0pt}
\setlength{\parskip}{3pt}
\setlength{\emergencystretch}{3em}
\widowpenalty=10000\clubpenalty=10000\displaywidowpenalty=10000
\setlist[enumerate]{leftmargin=*,itemsep=2pt,topsep=3pt}
\newcommand{\R}{\mathbb R}
\newcommand{\N}{\mathbb N}
\newcommand{\E}{\mathbb E}
\newcommand{\Pp}{\mathbb P}
\newcommand{\Q}{\mathbb Q}
\newcommand{\F}{\mathcal F}
\newcommand{\ind}{\mathbf 1}
\newcommand{\cE}{\mathcal E}
\newcommand{\Law}{\operatorname{Law}}
\newcommand{\sgn}{\operatorname{sgn}}
\newcommand{\Var}{\operatorname{Var}}
\newcommand{\dist}{\operatorname{dist}}
\newcommand{\KL}{\operatorname{KL}}
\newcommand{\TV}{\operatorname{TV}}
\newcommand{\Lip}{\operatorname{Lip}}
\DeclareMathOperator*{\esssup}{ess\,sup}
\DeclareMathOperator*{\argmax}{arg\,max}
\newtheorem{theorem}{Theorem}
\newtheorem{lemma}[theorem]{Lemma}
\newtheorem{proposition}[theorem]{Proposition}
\newtheorem{corollary}[theorem]{Corollary}
\theoremstyle{definition}
\newtheorem{example}[theorem]{Example}
\newcommand{\report}[3]{\clearpage\setcounter{theorem}{0}\renewcommand{\thetheorem}{#1.\arabic{theorem}}\renewcommand{\theHtheorem}{#1.\arabic{theorem}}\section*{#1: #2}\addcontentsline{toc}{section}{#1: #2}\textbf{Outcome:} #3.\par\textbf{Tools:} web browsing available; Python computation available. Literature checked on 9 October 2026.}
\newcommand{\stage}[2]{\subsection*{#1) #2}}
\newcommand{\unresolved}{\textbf{LABEL: UNRESOLVED.}}
\newcommand{\disproof}{\textbf{LABEL: FULL SOLUTION. SUBLABEL: COUNTEREXAMPLE/DISPROOF.}}

\begin{document}
\section*{C65-26: Critical exponential uniqueness for convex quadratic BSDEs}
\textbf{Outcome:} UNRESOLVED; subgradient densities and several uniqueness subclasses proved.\par\textbf{Tools:} web browsing available; Python computation available. Literature checked on 9 October 2026.
\subsection*{1) FORMAL RESTATEMENT}
\paragraph{Brownian conventions used in this report.}
Fix $T>0$ and an integer $d\ge1$. The probability space carries a $d$-dimensional Brownian motion $B$ with its completed, right-continuous natural filtration; $\F=\F_T$ and $\F_0$ is trivial modulo null sets. A scalar solution consists of a continuous adapted real process $Y$ and a predictable $\R^d$-valued process $Z$ such that
\[
 \int_0^T |Z_t|^2dt+\int_0^T|g(t,Y_t,Z_t)|dt<\infty\quad\Pp\text{-a.s.},
\]
and the displayed BSDE holds simultaneously for all times outside one null set. Equality of solutions means indistinguishability of $Y$ and $dt\otimes\Pp$-a.e. equality of $Z$. A process $X$ is of class (D) when $\{X_\tau:\tau\le T\text{ a stopping time}\}$ is uniformly integrable. Logs are natural, norms Euclidean, and $\sgn(0)=0$. At $T=0$ the assertions reduce to $Y_0=\xi$, with no nontrivial $Z$ component.

\paragraph{Precisely stated probabilistic tools.}
We use It\^o's formula and its Tanaka version for continuous semimartingales, Brownian martingale representation, and the following forms of change of measure. If $\theta$ is predictable and $\E\exp(\frac12\int_0^T|\theta|^2)<\infty$, Novikov's theorem says that
$D=\cE(\int\theta\,dB)$ is a uniformly integrable martingale; under $d\Q=D_Td\Pp$, $B-\int\theta dt$ is Brownian. Here $\cE(M)=\exp(M-\langle M\rangle/2)$. Also, for a continuous Brownian martingale, the BMO hypothesis
\[
 \esssup_{\tau\le T}\E\left[\left.\int_\tau^T|\theta_t|^2dt\right|\F_\tau\right]<\infty
\]
implies that $D$ is uniformly integrable and $D_T\in L^q$ for some $q>1$; this is the continuous-martingale BMO exponential theorem \cite{Kazamaki}. Each use below verifies the relevant integrability hypothesis. A local submartingale or supermartingale of class (D) is a true submartingale or supermartingale: apply the localized conditional inequalities and pass to the limit using uniform integrability. A local martingale of class (D) therefore equals the conditional expectation of its terminal value.

For every finite convex $g:\R^d\to\R$ with $0\le g(z)\le |z|^2/2$, and every $\F_T$-measurable $\xi$ with $\E e^{|\xi|}<\infty$, must any two solutions of
\[
 Y_t=\xi+\int_t^Tg(Z_s)ds-\int_t^TZ_s\cdot dB_s
\]
for which $e^{|Y|}$ is of class (D) coincide? This is an at-most-one assertion, not an assertion that every datum has a solution. The constants $1/2$ and $1$ are fixed. A stronger terminal exponential moment, or a positive lower curvature bound, may not be silently substituted.

The subdifferential is
$\partial g(z)=\{q:g(w)\ge g(z)+q\cdot(w-z)\text{ for every }w\}$.
The stress points are unbounded terminal data with exactly one exponential moment, flat regions of $g$, and change of measure at the boundary of entropy integrability.

\subsection*{2) QUICK LITERATURE/CONTEXT CHECK}
This is Problem 5.4 in Fan--Hu--Tang \cite{FHT}. Delbaen--Hu--Richou \cite{DHR} prove a critical-case result with strong convexity; their Proposition 3.1 already establishes the uniformly integrable subgradient density. We reprove the latter fact below in the present sign convention. Wang--Zhang--Fan \cite{WZF} impose additional structure and do not supply a general-convex endpoint theorem applicable here. These checks are not a claim of an exhaustive absence of later results.

\subsection*{3) ATTACK PLAN}
\textbf{Proof routes.} Use convex subgradients to change measure and compare solutions; alternatively transform entropic generators explicitly; finally identify what extra integrability makes the comparison rigorous. The first route is pursued up to its precise optional-sampling gap.

\textbf{Disproof routes.} Test zero and bounded terminal values, generators with flat cores, and terminal tails for which a change-of-measure density has infinite entropy. The first two do not produce a counterexample; the last falsifies a tempting auxiliary claim, not uniqueness itself.

\textbf{Landscape and tools.} This is critical stochastic comparison. The seven main tools are supporting hyperplanes (to bound subgradients), convex conjugacy (to interpret the bound), exponential It\^o transforms (to construct densities), Novikov after an intermediate change of measure (to avoid an unavailable moment), class-(D) localization (to justify sampling), conditional Jensen inequalities (for the entropic subclass), and entropy--exponential duality (to isolate sufficient extra integrability).

\subsection*{4) WORK}
\begin{lemma}[A sharp subgradient inequality]
For every $z\in\R^d$ and $q\in\partial g(z)$,
\[
 |q-z|^2\le |z|^2-2g(z),\qquad |q|\le2|z|.
\]
\end{lemma}
\begin{proof}
The supporting inequality at the point $w=q$ and the quadratic upper bound give
$g(z)+q\cdot(q-z)\le g(q)\le|q|^2/2$.
Rearranging proves the first assertion. Its right side is at most $|z|^2$, so the triangle inequality proves the second. Finite convex functions on $\R^d$ have nonempty subdifferentials. Their subdifferential graphs are closed, and here the values are compact by the second bound. A Borel measurable selection exists for a nonempty compact-valued Borel correspondence of this kind (the finite-dimensional measurable selection theorem). Composing such a selection with a predictable $Z$ gives a predictable $q_t\in\partial g(Z_t)$. The bound ensures $\int|q|^2<\infty$ a.s.
\end{proof}

\begin{proposition}[The critical subgradient density is a true density]
For any admissible solution $(Y,Z)$ and any predictable $q_t\in\partial g(Z_t)$, the stochastic exponential $\cE(\int q\,dB)$ is a uniformly integrable martingale.
\end{proposition}
\begin{proof}
Define the finite, increasing process
\[
 A_t=\int_0^t\bigl(|Z_s|^2/2-g(Z_s)\bigr)ds\ge0.
\]
The BSDE gives
\[
 D^0_t=\cE\Bigl(\int Z\,dB\Bigr)_t
       =\exp(Y_t-Y_0-A_t)\le e^{-Y_0}e^{Y_t}.
\]
The stopping family on the right is uniformly integrable, since it is dominated by $e^{-Y_0}e^{|Y_\tau|}$. Thus the positive local martingale $D^0$ is of class (D), is a true uniformly integrable martingale, and $\E D^0_T=1$. Its terminal value is strictly positive. Set $d\Q^0=D^0_Td\Pp$ and $B^0=B-\int Zdt$.

Put $r=q-Z$. The previous lemma yields $\frac12\int_0^T|r|^2dt\le A_T$. Moreover,
\[
 \E_{\Q^0}e^{A_T}
 =\E[D^0_Te^{A_T}]
 =e^{-Y_0}\E e^\xi<\infty.
\]
Novikov's condition therefore holds for $\int r\,dB^0$ under $\Q^0$. Its exponential $D^1$ is a uniformly integrable $\Q^0$-martingale of mean one. Direct addition of the logarithms gives
\[
 D^0_tD^1_t
 =\exp\left(\int_0^tq_s\,dB_s-\frac12\int_0^t|q_s|^2ds\right).
\]
Indeed the cross term $-\int r\cdot Z$ in $\int r\,dB^0$ completes the square. Consequently the terminal expectation of this positive local $\Pp$-martingale is
$\E[D^0_TD^1_T]=\E_{\Q^0}D^1_T=1$. A nonnegative local martingale is a supermartingale; equal expectations at zero and $T$ force equality of all its conditional supermartingale inequalities. It is thus the martingale closed by its integrable terminal value, hence uniformly integrable.
\end{proof}

This removes a frequently proposed but incorrect first gap: it is not necessary to prove Novikov directly for $q$ under $\Pp$.

\begin{theorem}[Uniqueness for bounded terminal data]
If $|\xi|\le K$ a.s., then the original solution class contains at most one solution, without strong convexity.
\end{theorem}
\begin{proof}
Since $g\ge0$, $Y$ is a local supermartingale. It is of class (D), because $|Y_\tau|\le e^{|Y_\tau|}$, so
$Y_t\ge\E[\xi\mid\F_t]\ge-K$.
It\^o's formula shows that $e^Y$ is a local submartingale, its drift being
$e^Y(|Z|^2/2-g(Z))dt\ge0$. Its class-(D) property gives
$e^{Y_t}\le\E[e^\xi\mid\F_t]\le e^K$.
Thus every admissible solution satisfies $|Y_t|\le K$.

Let $(\widetilde Y,\widetilde Z)$ be a second solution and choose $q\in\partial g(Z)$. The preceding proposition gives an equivalent probability $\Q$. With $\Delta Y=\widetilde Y-Y$ and $\Delta Z=\widetilde Z-Z$, Girsanov's theorem gives
\[
 d\Delta Y_t=
 \bigl[q_t\cdot\Delta Z_t-g(\widetilde Z_t)+g(Z_t)\bigr]dt
       +\Delta Z_t\,dB_t^{\Q}.
\]
The bracket is nonpositive by convexity. The process $\Delta Y$ is bounded, so localization makes it a true $\Q$-supermartingale. Since $\Delta Y_T=0$, it follows that $\Delta Y_t\ge0$. Reversing the two solutions gives the reverse inequality. Equality first at every rational time, then continuity, gives indistinguishability. The difference equation now has zero quadratic variation, so $\int_0^T|\Delta Z|^2dt=0$ a.s.
\end{proof}

\begin{theorem}[The entire entropic ray is solved]
For $g(z)=\lambda|z|^2/2$, $0\le\lambda\le1$, the original question has a unique solution for every allowed terminal value. For $\lambda>0$ it is
\[
 Y_t=\lambda^{-1}\log\E[e^{\lambda\xi}\mid\F_t],
\]
and for $\lambda=0$ it is $Y_t=\E[\xi\mid\F_t]$.
\end{theorem}
\begin{proof}
For an admissible solution and $\lambda>0$, It\^o's formula makes $e^{\lambda Y}$ a local martingale. It is dominated in absolute value by $1+e^{|Y|}$, and hence is of class (D). Its terminal conditional expectation proves uniqueness of the displayed $Y$.

Conversely let $M_t=\E[e^{\lambda\xi}\mid\F_t]$. Brownian representation supplies $dM=H\,dB$. The continuous martingale $M$ is strictly positive throughout $[0,T]$ a.s.: if it hit zero, optional sampling of the nonnegative martingale would force its strictly positive terminal value to vanish on that hitting event. A continuous strictly positive path on $[0,T]$ has a positive minimum. Local martingale representation and localization therefore make $Z=H/(\lambda M)$ square-integrable in time a.s. It\^o's logarithm formula verifies the BSDE.

Conditional power-mean and Jensen inequalities imply
\[
 e^{Y_t}\le\E[e^\xi\mid\F_t],\qquad
 e^{-Y_t}\le e^{-\E[\xi\mid\F_t]}\le\E[e^{-\xi}\mid\F_t].
\]
The first inequality is conditional concavity of $x\mapsto x^\lambda$ on $(0,\infty)$; the second uses
$\lambda^{-1}\log\E e^{\lambda\xi}\ge\E\xi$ conditionally. Therefore $e^{|Y|}\le e^Y+e^{-Y}$ is dominated by the sum of two uniformly integrable closed martingales and is of class (D). The case $\lambda=0$ follows by martingale representation and conditional Jensen directly. In each case the $Z$ component is unique by quadratic variation.
\end{proof}

\begin{proposition}[A sufficient, but nonautomatic, entropy condition]
For two admissible solutions, suppose a subgradient density can be chosen for each solution with finite entropy $\E[D_T\log D_T]<\infty$. Then the two solutions coincide.
\end{proposition}
\begin{proof}
Fix one such density $D_T$. Young's inequality $ab\le a\log a-a+e^b$, for $a>0,b\ge0$, gives
\[
 D_T|Y_\tau|\le D_T\log D_T-D_T+e^{|Y_\tau|}.
\]
The first two terms form a single integrable variable (the negative part of $a\log a$ is bounded), and the last family is uniformly integrable. Thus the family $D_T|Y_\tau|$ is uniformly integrable under $\Pp$. Also $\sup_\tau\E|Y_\tau|<\infty$, so $\Pp(|Y_\tau|>K)\to0$ uniformly in $\tau$. Uniform absolute continuity of the preceding integrals proves
$\sup_\tau\E[D_T|Y_\tau|\ind_{\{|Y_\tau|>K\}}]\to0$.
Hence $Y$ is of class (D) under $\Q$. This argument also applies to $\widetilde Y$ using the same density. The convexity comparison from the bounded theorem is therefore valid with class (D) in place of boundedness. It gives one order; the density chosen for the other solution gives the other order.
\end{proof}

\begin{example}[Infinite entropy really can occur]
Let $U=\Phi(B_T^{(1)}/\sqrt T)$, where $\Phi$ is the standard normal distribution function, so $U$ is uniform on $(0,1)$. Set
\[
 \xi=\begin{cases}
 \displaystyle\log\frac{1}{U(\log(1/U))^2},&U<e^{-4},\\
 0,&U\ge e^{-4}.
 \end{cases}
\]
Then $\xi\ge0$ and, by $r=\log(1/U)$,
\[
 \E e^{|\xi|}=1-e^{-4}+\int_4^\infty r^{-2}dr
 =1-e^{-4}+\tfrac14<\infty,
\]
whereas
\[
 \E[e^\xi\xi]=\int_4^\infty\frac{r-2\log r}{r^2}dr=\infty.
\]
For $g(z)=|z|^2/2$, the solved entropic formula gives the subgradient density
$D_T=e^\xi/\E e^\xi$, whose entropy is infinite. This example belongs to a uniqueness subclass; it refutes automatic finite entropy, not the original uniqueness statement.
\end{example}

\subsection*{5) VERIFICATION}
\textbf{Tiny cases and computation.} Constant terminal values are covered by bounded-data uniqueness. Both $g=0$ and $g=|z|^2/2$ are covered at the exact endpoint. The script checks the subgradient inequality for the flat-core convex generator
$g(z)=\frac12((|z|-1)^+)^2$ on 20,001 points; its minimum slack is zero. Truncated entropy integrals at upper limits $10,100,10^4$ are about $0.384,2.138,6.633$, consistent with the analytically proved divergence.

\textbf{Adversarial review.} Uniform integrability of a density alone does not transfer class (D) to the changed probability. The first proposition proves only that the probability exists; it does not prove the comparison step for unbounded differences. No positive curvature is used in the bounded-data theorem. All stochastic integrands remain pathwise square-integrable after equivalent changes of measure. Convex subgradient selection, the signs of the drift, and both comparison directions are explicitly accounted for.

\subsection*{6) FINAL}
\unresolved

\textbf{Strongest checked partial result.} Every admissible solution has a uniformly integrable subgradient exponential; uniqueness holds for bounded terminal data and, for arbitrary permitted terminal data, for every $g(z)=\lambda|z|^2/2$, $0\le\lambda\le1$. Finite entropy for comparison densities is another proved sufficient condition.

\textbf{Exact first remaining gap in the selected route.} For two arbitrary admissible solutions, prove that the convexity-induced local supermartingale $\widetilde Y-Y$, under the first solution's subgradient probability, satisfies the terminal comparison inequality $\widetilde Y_t-Y_t\ge\E_{\Q}[\widetilde Y_T-Y_T\mid\F_t]=0$, without imposing a stronger moment or curvature hypothesis. It would suffice to establish the required one-sided uniform integrability; it has not been established here.

\textbf{Next three moves.} (1) Prove a one-sided sampling estimate exploiting the Bregman divergence $g(\widetilde Z)-g(Z)-q\cdot(\widetilde Z-Z)$, rather than entropy alone. (2) Decompose directions with positive and zero curvature and test whether entropic comparison can be applied only in the former. (3) Investigate flat-core generators with explicit borderline terminal tails by stopped, finite-horizon approximations, checking the exact exponential class-(D) condition.

\textbf{Likely structure of a counterexample, if one exists.} It must escape bounded terminal data and the quadratic-ray generators. A plausible candidate would combine unbounded critical tails, nonuniform curvature, and loss of the particular changed-measure sampling property. The infinite-entropy example alone is not such a counterexample.

\clearpage
\begin{thebibliography}{99}
\bibitem{Kazamaki}
Norihiko Kazamaki. \emph{Continuous Exponential Martingales and BMO}. Lecture Notes in Mathematics 1579, Springer, 1994. DOI: 10.1007/BFb0073585. Continuous BMO stochastic exponentials and reverse H\"older inequalities. \url{https://doi.org/10.1007/BFb0073585}.

\bibitem{FHT}
Shengjun Fan, Ying Hu and Shanjian Tang. \emph{1D nonlinear backward stochastic differential equations: a unified theory and applications}. arXiv:2309.06233v2, 11 February 2026. In particular Problems 5.1--5.5, PDF pp.\ 32--33; for C65-45 see (A1), (4.13)--(4.14), PDF p.\ 31. \url{https://arxiv.org/pdf/2309.06233v2}.

\bibitem{DHR}
Freddy Delbaen, Ying Hu and Adrien Richou. \emph{On the uniqueness of solutions to quadratic BSDEs with convex generators and unbounded terminal conditions: The critical case}. Discrete and Continuous Dynamical Systems 35(11) (2015), 5273--5283. DOI: 10.3934/dcds.2015.35.5273. Proposition 3.1 and Theorem 4.1. Preprint: \url{https://arxiv.org/abs/1303.4859}; publisher: \url{https://www.aimsciences.org/article/doi/10.3934/dcds.2015.35.5273}.

\bibitem{WZF}
Yan Wang, Yaqi Zhang and Shengjun Fan. \emph{On the uniqueness of solutions to quadratic BSDEs with non-convex generators and unbounded terminal conditions: the certain exponential moment case}. arXiv:2409.13463v2, 3 August 2026. \url{https://arxiv.org/abs/2409.13463v2}.

\end{thebibliography}
\end{document}
